\mwnewpart
\renewcommand\runningtitle{Section 6. Commands}
\chapter{Commands}
\startcontents[sec]

This Section describes every RetrOS/LT command, in alphabetical order.
\term{Internal} commands are part of \file{COMMAND.COM} and are always
available. \term{External} commands are files on the DOS diskette, or in
the directory Install put them in; that directory must be the current one
or be named in the \typed{PATH}.

In the formats, \var{d:} is a drive, \var{path} a directory path,
\var{file} a file name, which may usually have a drive and path before it
and wildcards in it. Items in square brackets may be left out.

The commands for batch files are \typed{ECHO}, \typed{FOR}, \typed{GOTO},
\typed{IF}, \typed{PAUSE}, \typed{REM} and \typed{SHIFT}. \typed{BASIC},
\typed{DEBUG}, \typed{LINK} and \typed{EXE2BIN} are on the Supplemental
Programs diskette.

\cmd{ASSIGN}{External}{Sends requests for one drive to another.}
\fmt{ASSIGN [\var{x}=\var{y} [...]]}
\typed{ASSIGN A=C} sends requests for drive A to drive C, for programs that
look for their files only on drive A. \typed{ASSIGN} alone puts every drive
back. Do not use \typed{FORMAT}, \typed{DISKCOPY}, \typed{DISKCOMP} or
\typed{PRINT} while drives are assigned. \typed{ASSIGN} stays in memory once
it has been used.

\cmd{BACKUP}{External}{Copies files from a disk to diskettes for
safekeeping, across as many diskettes as needed.}
\fmt{BACKUP \var{d:}[\var{path}][\var{file}] \var{d:} [/S] [/M] [/A] [/D:\var{mm-dd-yy}]}
\begin{opts}
\opt{/S}{Includes the directories below the source.}
\opt{/M}{Only files changed since they were last backed up.}
\opt{/A}{Adds to the files already on the backup diskettes, beginning
  with the last one, instead of erasing them.}
\opt{/D:}{Only files changed on or after the date.}
\end{opts}
The target diskettes must be formatted. Only \typed{RESTORE} can read what
\typed{BACKUP} writes. \typed{BACKUP} ends with exit code 0 if it backed up
the files, 1 if there were none to back up, 3 if it was stopped with
\keys{Ctrl}{Break}, and 4 after an error; \typed{IF ERRORLEVEL} can test
it. See Section~4.

\cmd{BASIC, BASICA}{External}{Starts BASIC.}\ixs{BASIC}
\fmt{BASICA [\var{file}] [/F:\var{files}] [/S:\var{size}] [/C:\var{size}] [/M:\var{memory}]}
\fmt{BASIC /ROM}
\begin{opts}
\opt{\var{file}}{A BASIC program to load and run; \file{.BAS} is assumed.}
\opt{/F:, /S:}{IBM's number of files and largest random record. They are
  accepted, and not needed: this BASIC makes room for files and records
  as they are opened.}
\opt{/C:}{Bytes for receiving from the serial ports, half for each; 512
  unless given.}
\opt{/M:}{The highest memory location BASIC is to use, within its own
  64 KB, leaving the rest for machine language routines.}
\opt{/ROM}{Starts Cassette BASIC, in the computer's read-only memory,
  instead. There is no coming back but restarting the computer.}
\end{opts}
\file{BASICA.EXE} and \file{BASIC.COM} are on the Supplemental Programs
diskette. \typed{SYSTEM} leaves BASIC. See ``BASIC'' in Section~3.

\cmd{BREAK}{Internal}{Sets how often RetrOS/LT checks for
\keys{Ctrl}{Break}.}
\fmt{BREAK [ON|OFF]}
\typed{BREAK} alone shows the setting. See \typed{BREAK=} in Section~5.

\cmd{CHDIR (CD)}{Internal}{Shows or changes the current directory.}
\fmt{CD [\var{d:}][\var{path}]}
\typed{CD} alone shows the current directory. \typed{CD ..} goes up one
directory; \typed{CD \bs} goes to the root. Naming a drive changes the
current directory of that drive without making it the current drive.

\cmd{CHKDSK}{External}{Checks a disk and reports its space and the
computer's memory.}
\fmt{CHKDSK [\var{d:}][\var{file}] [/F] [/V]}
\begin{opts}
\opt{/F}{Fixes the errors it finds, after asking.}
\opt{/V}{Shows the name of every file as it is checked.}
\end{opts}
With a file name, \typed{CHKDSK} also reports whether those files are in
one piece. Do not use it on an assigned drive. See Section~4.

\cmd{CLS}{Internal}{Clears the screen.}
\fmt{CLS}

\cmd{COMMAND}{External}{Starts another copy of the command processor.}
\fmt{COMMAND [\var{d:}][\var{path}] [/P] [/C \var{command}]}
\begin{opts}
\opt{\var{path}}{Where \file{COMMAND.COM} is to be found when it must be
  read again.}
\opt{/P}{Makes the new copy permanent, and runs \file{AUTOEXEC.BAT}.}
\opt{/C \var{command}}{Runs the command and then returns.}
\end{opts}
\typed{EXIT} ends a copy that is not permanent.

\cmd{COMP}{External}{Compares two files, or two sets of files, byte for
byte.}
\fmt{COMP [\var{file1}] [\var{file2}]}
If you leave out a name, \typed{COMP} asks for it. The second may be a
drive or directory only, or have wildcards, and then stands for files of
the same names. \typed{COMP} reports each difference by its offset in the
file, in hexadecimal, with the two bytes, and gives up on a pair after ten
differences. Files of different sizes are not compared. At the end it asks
\shown{Compare more files (Y/N)?}.

\cmd{COPY}{Internal}{Copies files, or joins them.}
\fmt{COPY [/A|/B] \var{source}[/A|/B] [\var{target}[/A|/B]] [/V]}
\fmt{COPY \var{source1}+\var{source2}[+...] [\var{target}] [/V]}
\begin{opts}
\opt{/A}{Treats the file as text, ending at the first end-of-file mark
  (\keys{Ctrl}{Z}).}
\opt{/B}{Copies the whole file, whatever it holds.}
\opt{/V}{Verifies that what was written can be read.}
\end{opts}
If the target is left out, the copy goes in the current directory under
the same name. The source or target may be a device: \typed{COPY CON
NOTE.TXT} copies what you type into a file, until you press \key{F6} and
\key{Enter}; \typed{COPY NOTE.TXT PRN} prints it. With \typed{+} between
them, the sources are joined into the target. \typed{COPY} replaces an
existing target without asking.

\cmd{CTTY}{Internal}{Changes the device RetrOS/LT takes commands from.}
\fmt{CTTY \var{device}}
\typed{CTTY COM1} lets a terminal on the serial port run RetrOS/LT;
\typed{CTTY CON} from that terminal gives the keyboard and screen back.

\cmd{DATE}{Internal}{Shows or sets the date.}
\fmt{DATE [\var{mm-dd-yyyy}]}
Without a date, \typed{DATE} shows the date and asks for a new one; press
\key{Enter} to keep it. Give the year in four digits. RetrOS/LT keeps the
date only while the computer is on.

\cmd{DEBUG}{External}{A program for examining and changing memory, files
and disk sectors, and for testing programs.}
\fmt{DEBUG [\var{file} [\var{parameters}]]}
Appendix~C lists its commands. \typed{Q} leaves it.

\cmd{DEL (ERASE)}{Internal}{Erases files.}
\fmt{DEL \var{file}}
\typed{DEL *.*} asks \shown{Are you sure (Y/N)?} first; otherwise
\typed{DEL} does not ask.

\cmd{DIR}{Internal}{Lists the files in a directory.}
\fmt{DIR [\var{d:}][\var{path}][\var{file}] [/P] [/W]}
\begin{opts}
\opt{/P}{Stops after each screen.}
\opt{/W}{Shows the names only, five to a line.}
\end{opts}
\typed{DIR} shows each file's name, size, and the date and time it was last
changed, and the space left on the disk. It does not show the hidden system
files.

\cmd{DISKCOMP}{External}{Compares two diskettes, track for track.}
\fmt{DISKCOMP [\var{d:} [\var{d:}]] [/1] [/8]}
\typed{/1} compares only the first side, and \typed{/8} only 8 sectors of
each track. \typed{DISKCOMP} reports each track that differs, by side and
track; \shown{Diskettes compare ok} if none does.

\cmd{DISKCOPY}{External}{Copies a whole diskette to another of the same
kind.}
\fmt{DISKCOPY [\var{d:} [\var{d:}]] [/1]}
\typed{DISKCOPY} formats the target as it copies if it is not formatted
the same way. \typed{/1} copies only the first side. See Section~4.

\cmd{ECHO}{Internal, batch}{Shows a message, or turns the showing of
batch commands on or off.}
\fmt{ECHO [ON|OFF|\var{message}]}
\typed{ECHO} alone shows the setting.

\cmd{EDLIN}{External}{A line editor for text files.}
\fmt{EDLIN \var{file} [/B]}
\typed{/B} reads the whole file, past any \keys{Ctrl}{Z}. Appendix~C lists
its commands. \typed{E} saves the file and leaves; \typed{Q} leaves without
saving.

\cmd{EXE2BIN}{External}{Converts a program from \file{.EXE} form to
\file{.COM} form, or to a memory image.}
\fmt{EXE2BIN \var{file} [\var{target}]}
The program must have no stack segment and be no larger than 64 KB.

\cmd{EXIT}{Internal}{Leaves a copy of the command processor started by
\typed{COMMAND} or by a program.}
\fmt{EXIT}

\cmd{FC}{External}{Compares two files and shows where they differ.}
\fmt{FC [/\var{n}] [/B] [/C] [/W] \var{file1} \var{file2}}
\begin{opts}
\opt{/B}{Compares byte for byte, as \typed{COMP} does, for programs.}
\opt{/C}{Ignores the difference between capitals and small letters.}
\opt{/W}{Treats spaces and tabs as a single space, and ignores them at the
  ends of lines.}
\opt{/\var{n}}{The number of lines, 1 to 9, that must match again before
  a difference is taken to have ended.}
\end{opts}
For text, \typed{FC} shows the lines that differ, with a line either side.
\typed{FC} is not in PC DOS 2.10; it is Microsoft's, from the same
source.

\cmd{FDISK}{External}{Sets up the partitions of a fixed disk.}
\fmt{FDISK}
See ``Partitioning a Fixed Disk'' in Section~4.

\cmd{FIND}{External}{Shows the lines of a file that contain a string of
characters.}
\fmt{FIND [/V] [/C] [/N] "\var{string}" [\var{file} ...]}
\begin{opts}
\opt{/V}{Shows the lines that do not contain it.}
\opt{/C}{Shows only how many lines contain it.}
\opt{/N}{Shows each line's number.}
\end{opts}
Capitals and small letters differ. Without a file, \typed{FIND} reads what
is piped to it: \typed{DIR | FIND "BAT"}.

\cmd{FOR}{Internal, batch}{Runs a command once for each of a set of files
or words.}
\fmt{FOR \%\%\var{v} IN (\var{set}) DO \var{command}}
In a batch file the variable has two percent signs, \typed{\%\%F}; typed at
the prompt, one: \typed{FOR \%F IN (*.TXT) DO TYPE \%F}.

\cmd{FORMAT}{External}{Prepares a diskette or a fixed disk to hold files.}
\fmt{FORMAT \var{d:} [/S] [/V] [/1] [/8] [/B]}
\begin{opts}
\opt{/S}{Copies the system files and \file{COMMAND.COM} to it, so that
  it starts the computer.}
\opt{/V}{Asks for a volume label.}
\opt{/1}{Formats one side only.}
\opt{/8}{8 sectors a track.}
\opt{/B}{8 sectors a track, leaving room for \typed{SYS} to put the system
  on later.}
\end{opts}
See Section~4.

\cmd{GOTO}{Internal, batch}{Jumps to a label in a batch file.}
\fmt{GOTO \var{label}}
A label is a line that begins with a colon: \typed{:AGAIN}.

\cmd{GRAPHICS}{External}{Lets \keys{Shift}{PrtSc} print a graphics screen.}
\fmt{GRAPHICS}
After \typed{GRAPHICS}, \keys{Shift}{PrtSc} prints the color graphics
screen on the IBM Graphics Printer, turned on its side: in 640 by 200 mode
dot for dot, and in 320 by 200 mode with each dot made larger and the four
colors printed as four shades. Text screens are printed as before.
\typed{GRAPHICS} stays in memory.

\cmd{IF}{Internal, batch}{Runs a command only if a condition holds.}
\fmt{IF [NOT] ERRORLEVEL \var{n} \var{command}}
\fmt{IF [NOT] \var{string1}==\var{string2} \var{command}}
\fmt{IF [NOT] EXIST \var{file} \var{command}}
\typed{ERRORLEVEL \var{n}} is true if the last program's exit code was
\var{n} or more. \typed{BACKUP} and \typed{RESTORE} set one.

\cmd{INSTALL}{External}{Puts RetrOS/LT on the fixed disk.}
\fmt{INSTALL}
Run from the DOS diskette. See Section~2.

\cmd{LINK}{External}{Links object modules made by an assembler or a
compiler into a program.}
\fmt{LINK [\var{objects}[,\var{program}[,\var{map}[,\var{libraries}]]]][;]}
\typed{LINK} asks for whatever is left out; a semicolon takes the usual
answers for the rest. It is Microsoft's linker, version 2.00.

\cmd{MKDIR (MD)}{Internal}{Makes a directory.}
\fmt{MD [\var{d:}]\var{path}}

\cmd{MODE}{External}{Sets up the printers, the serial port and the
display.}
\fmt{MODE LPT\var{n}:[\var{chars}][,[\var{lines}][,P]]}
\fmt{MODE LPT\var{n}:=COM1}
\fmt{MODE COM1:\var{baud}[,\var{parity}[,\var{data}[,\var{stop}[,P]]]]}
\fmt{MODE \var{display}}
\fmt{MODE [\var{display}],R|L[,T]}
\begin{longopts}
\lopt{LPT\var{n}:\var{chars},\var{lines}}{80 or 132 characters a line, and
  6 or 8 lines an inch, on an IBM printer. \typed{P} keeps trying when the
  printer is not ready, instead of reporting an error.}
\lopt{LPT\var{n}:=COM1}{Sends what is printed on that printer to the
  serial port. \typed{MODE LPT\var{n}:} alone stops it.}
\lopt{COM1:\var{baud}}{110, 150, 300, 600, 1200, 2400, 4800 or 9600 (the
  first two digits are enough); parity \typed{N}, \typed{O} or \typed{E}
  (\typed{E} if not given); 7 or 8 data bits (7); 1 or 2 stop bits (1, or 2
  at 110 baud). \typed{P} keeps trying when the port is not ready, for a
  serial printer.}
\lopt{\var{display}}{\typed{40} or \typed{80} columns; \typed{BW40},
  \typed{BW80}, \typed{CO40} or \typed{CO80} on the color display; or
  \typed{MONO} for the monochrome display.}
\lopt{R, L, T}{Moves the picture on the color display right or left, two
  characters in 80 columns and one in 40. \typed{T} shows a test pattern
  and asks whether it is lined up.}
\end{longopts}
\typed{MODE} stays in memory when it reroutes or retries a printer or
shifts the display.

\cmd{MORE}{External}{Shows what is piped to it one screen at a time.}
\fmt{MORE < \var{file}}
\fmt{\var{command} | MORE}
After each screen it shows \shown{-- More --}; press a key for the next.

\cmd{PATH}{Internal}{Sets the directories searched for programs and batch
files.}
\fmt{PATH [[\var{d:}]\var{path}[;[\var{d:}]\var{path}...]]}
\fmt{PATH ;}
\typed{PATH} alone shows the path; \typed{PATH ;} removes it.

\cmd{PAUSE}{Internal, batch}{Stops a batch file until a key is pressed.}
\fmt{PAUSE [\var{message}]}
It shows the message, if echoing is on, and \shown{Strike a key when
ready . . .}.

\cmd{PRINT}{External}{Prints text files while you go on working.}
\fmt{PRINT [\var{file} ...] [/T] [/C] [/P]}
\begin{opts}
\opt{/T}{Cancels every file in the queue.}
\opt{/C}{Cancels the files named after it.}
\opt{/P}{Adds the files named after it (the usual action).}
\end{opts}
The first time it is run, \typed{PRINT} asks \shown{Name of list device
[PRN]:}; press \key{Enter} for the printer, or name another device. It then
stays in memory. The queue holds ten files. \typed{PRINT} alone shows the
queue.

\cmd{PROMPT}{Internal}{Sets what the prompt shows.}
\fmt{PROMPT [\var{text}]}
See ``The Environment'' in Section~5 for the codes. \typed{PROMPT \$p\$g}
shows the current directory.

\cmd{RECOVER}{External}{Rescues what can be read of a file that has a bad
sector, or of every file on a disk whose directory is damaged.}
\fmt{RECOVER \var{file}}
\fmt{RECOVER \var{d:}}
With a file name, \typed{RECOVER} keeps what can be read and marks the bad
sectors. With a drive alone, it makes every file into one named
\file{FILE\var{nnnn}.REC} in the root directory; use it only when the
directory itself cannot be read.

\cmd{REM}{Internal, batch}{A remark, ignored by RetrOS/LT.}
\fmt{REM [\var{text}]}

\cmd{RENAME (REN)}{Internal}{Changes a file's name.}
\fmt{REN \var{file} \var{newname}}
The new name may not have a drive or path; the file stays where it is.

\cmd{RESTORE}{External}{Puts back files saved by \typed{BACKUP}.}
\fmt{RESTORE \var{d:} [\var{d:}][\var{path}]\var{file} [/S] [/P]}
\begin{opts}
\opt{/S}{Includes the directories below.}
\opt{/P}{Asks before replacing a file that is read-only or that has been
  changed since it was backed up.}
\end{opts}
The first drive is the one with the backup diskettes. Files are restored to
the directories they came from, which are made if need be. Exit codes are
as for \typed{BACKUP}.

\cmd{RMDIR (RD)}{Internal}{Removes a directory.}
\fmt{RD [\var{d:}]\var{path}}
The directory must be empty, and not the current one.

\cmd{SET}{Internal}{Shows, sets or removes a value in the environment.}
\fmt{SET [\var{name}=[\var{value}]]}
\typed{SET} alone shows them all. With nothing after the equal sign, the
name is removed.

\cmd{SHIFT}{Internal, batch}{Moves each batch parameter down one place,
so that a batch file can use more than nine.}
\fmt{SHIFT}

\cmd{SORT}{External}{Sorts lines of text.}
\fmt{SORT [/R] [/+\var{n}]}
\begin{opts}
\opt{/R}{Sorts in reverse order.}
\opt{/+\var{n}}{Sorts on the characters from column \var{n}.}
\end{opts}
\typed{SORT} reads what is piped or redirected to it: \typed{SORT <
NAMES.TXT > SORTED.TXT}.

\cmd{SYS}{External}{Copies the system files to a formatted disk.}
\fmt{SYS \var{d:}}
See ``Making a Disk Start the Computer'' in Section~4.

\cmd{TIME}{Internal}{Shows or sets the time.}
\fmt{TIME [\var{hh}:\var{mm}[:\var{ss}[.\var{xx}]]]}
Use the 24-hour clock: \typed{TIME 14:30}.

\cmd{TREE}{External}{Shows the directories of a disk.}
\fmt{TREE [\var{d:}] [/F]}
\typed{TREE} lists each directory with the directories in it.
\typed{/F} lists the files in each as well. Add \typed{| MORE} for a large
disk.

\cmd{TYPE}{Internal}{Shows a file on the screen.}
\fmt{TYPE \var{file}}

\cmd{VER}{Internal}{Shows the version of RetrOS/LT.}
\fmt{VER}
\begin{dialogue}
A>\textbf{VER}

RetrOS/LT Version 1.0, DOS 2.10
\end{dialogue}

\cmd{VERIFY}{Internal}{Sets whether what is written to a disk is read back
and checked.}
\fmt{VERIFY [ON|OFF]}
Verifying makes writing slower. \typed{VERIFY} alone shows the setting.

\cmd{VOL}{Internal}{Shows a disk's volume label.}
\fmt{VOL [\var{d:}]}

\stopcontents[sec]
